\renewcommand{\bibname}{Quellen und Bandzuordnung}
\begin{thebibliography}{9}
\addcontentsline{toc}{chapter}{Quellen und Bandzuordnung}
\footnotesize
\setlength{\itemsep}{6pt}
\setlength{\parskip}{0pt}
\interlinepenalty=10000

\bibitem{project}
Martin Kunze: \emph{Die Grundlagen der Mathematik}.
Band~15: Begriff der totalen Ordnung; minimale und maximale
Elemente; Existenz und Definition der binären Minimums-
und Maximumswerte. Diese Kernaussagen werden in den
Kapiteln~\ref{b15guide:ch:foundations} und
\ref{b15guide:ch:extrema} zusammengefaßt.
Die dort ebenfalls bewiesenen allgemeinen Rechengesetze
und die Einordnung der natürlichen Zahlen werden nicht wiederholt.
Die übrigen Kapitel ergänzen elementare Standardfolgen der
vollständigen Vergleichbarkeit.
Format nach den Leitfäden zu den Bänden~3, 4, 5, 28 und~37.

\bibitem{linear}
The mathlib Community: \emph{Mathlib.Order.Defs.LinearOrder},
offizielle Dokumentation, abgerufen am 23.~Oktober 2026.
Abgleich von Trichotomie und Negationsregeln über
\path{lt_trichotomy}, \path{not_le}, \path{not_lt};
von Auswahlwerten und disjunktiven Vergleichsregeln über
\path{min_def}, \path{max_def}, \path{min_le_iff},
\path{le_max_iff}, \path{min_lt_iff}, \path{lt_max_iff}.
\url{https://leanprover-community.github.io/mathlib4_docs/Mathlib/Order/Defs/LinearOrder.html}.

\bibitem{intervals}
The mathlib Community: \emph{Mathlib.Order.Interval.Set.LinearOrder},
offizielle Dokumentation, abgerufen am 23.~Oktober 2026.
Vergleichsstellen für komplementäre Strahlen und ihre Inklusionskriterien:
\path{Set.compl_Iic}, \path{Set.compl_Iio},
\path{Set.Iio_subset_Iio_iff}, \path{Set.Ioi_subset_Ioi_iff}.
\url{https://leanprover-community.github.io/mathlib4_docs/Mathlib/Order/Interval/Set/LinearOrder.html}.

\bibitem{convex}
The mathlib Community: \emph{Mathlib.Order.Interval.Set.OrdConnected},
offizielle Dokumentation, abgerufen am 23.~Oktober 2026.
Begriffsabgleich der Ordnungskonvexität und der unorientierten
Intervalle über \path{Set.OrdConnected} und
\path{Set.ordConnected_iff_uIcc_subset}.
\url{https://leanprover-community.github.io/mathlib4_docs/Mathlib/Order/Interval/Set/OrdConnected.html}.

\bibitem{sections}
The mathlib Community: \emph{Mathlib.Order.UpperLower.Basic},
offizielle Dokumentation, abgerufen am 23.~Oktober 2026.
Vergleichbarkeit von Anfangs- und Endabschnitten unter Inklusion:
\path{IsLowerSet.total} und \path{IsUpperSet.total}.
\url{https://leanprover-community.github.io/mathlib4_docs/Mathlib/Order/UpperLower/Basic.html}.
\bibitem{cover}
The mathlib Community: \emph{Mathlib.Order.Cover},
offizielle Dokumentation, abgerufen am 23.~Oktober 2026.
Begriffsabgleich der unmittelbaren Nachbarschaft über \path{CovBy};
der Leitfaden verwendet ausschließlich den Spezialfall totaler Ordnungen.
\url{https://leanprover-community.github.io/mathlib4_docs/Mathlib/Order/Cover.html}.

\bibitem{lex}
The mathlib Community: \emph{Mathlib.Order.Prod.Lex.Basic},
offizielle Dokumentation, abgerufen am 23.~Oktober 2026.
Vergleichsquelle zur lexikographischen Ordnung auf einem
Produkt total geordneter Träger.
\url{https://leanprover-community.github.io/mathlib4_docs/Mathlib/Order/Prod/Lex/Basic.html}.

\end{thebibliography}

\begingroup
\footnotesize
\noindent\textbf{Umfang des Quellenabgleichs.}
Die externen Quellen dienen dem Abgleich der angegebenen Begriffe
und ausgewählten Standardregeln. Sie werden nicht als wortgleicher
Nachweis jeder Formulierung angeführt.
Insbesondere die zusammengefaßten Charakterisierungen, die
Schnittfallunterscheidung und die endliche Helly-Eigenschaft sind
hier in der gemeinsamen Mengensprache des Leitfadens formuliert.
Die Definition eines Schnitts ist ausdrücklich in
Definition~\ref{b15guide:def:cut} festgelegt.
Es werden weder neue Forschungsergebnisse noch historische
Erstentdeckungen beansprucht. Die Aussagen wurden redaktionell
auf ihre Voraussetzungen geprüft; eine maschinelle
Formalverifikation dieses Leitfadens wird nicht behauptet.
\endgroup
